% Cote 119: offprints, reports, courses, CV and letters (typed and printed part): archive cote 119, site page 291 (PDF page 292).
% typescript. Draft compared: none.
% AI-transcribed by Claude Opus 5.5 (claude-opus-5-5, adaptive thinking, effort high), Message Batches API, batch r119_full (msgbatch_01DwZNm5Y7YsKBWyAN4sFs8w), 2026-11-13; not yet checked by eye.
\scanpage{119/291}
\mspage{4}

Les conjectures de Milnor (\uline{Invent}.\uline{Math}. 9, 1970) disent que ces homomor-%
phismes sont des \uline{isomorphismes}. Cela a été démontré pour $n < 4$ (Merkurjev-%
Suslin, Arason, Rost) et il y a des résultats partiels pour $n$ \hand{$\geqslant$} 4.
Le cours s'est borné à citer ces énoncés sans en donner de démonstrations.
Il a été complété par deux exposés de B.Kahn sur les formes de Pfister
et leurs invariants cohomologiques.

§3. Applications et exemples

3.1. Invariants à valeurs dans $H^3(k, Z/2Z)$ : le cas du groupe spinoriel

Soit $q$ une forme quadratique non dégénérée sur $k$, et soit $x$ un
élément de $H^1(k, \mathrm{Spin}(q))$. Si l'on tord $q$ par $x$, on obtient une forme
quadratique $q_x$ de même rang que $q$. D'après (2.2.1) et (2.2.2), les
invariants $w_1$ et $w_2$ de $q_x$ sont les mêmes que ceux de $q$. Il en résulte que
l'élément $q_x - q$ de l'anneau de Witt $W_k$ appartient au cube $I_k^3$ de
l'idéal d'augmentation $I_k$ .En utilisant l'homomorphisme
\[ I_k^3/I_k^4 \qquad \to \qquad H^3(k, Z/2Z) \]
construit par Arason (qui est en fait un isomorphisme, cf. n° 2.3), on
obtient un élément de $H^3(k, Z/2Z)$ que nous noterons $i(x)$. On a :
\[ (3.1.1) \qquad i(x) = 0 \qquad \hand{$\Longleftrightarrow$} \qquad q_x \hand{$\equiv$} q \quad (\mathrm{mod}\ I_k^4). \]
On a ainsi défini une application canonique
\[ (3.1.2) \qquad \hand{$i$} : \quad H^1(k, \mathrm{Spin}(q)) \quad \to \ H^3(k, Z/2Z). \]

3.2. Invariants à valeurs dans $H^3(k, Z/2Z)$ : cas général

Prenons pour $G$ un groupe semi-simple \uline{simplement connexe} déployé,
et choisissons une représentation irréductible \hand{$\rho$} de $G$ dans un
espace vectoriel $V$ de dimension $n$. Supposons \hand{$\rho$} orthogonale, ce qui
est par exemple le cas si $G$ est de l'un des types $G_2$ , $F_4$ ou $E_8$ . Il existe
alors une forme quadratique non dégénérée $q$ sur $V$ qui est invariante
par \hand{$\rho$}$(G)$. On obtient ainsi un homomorphisme $G \to O(q)$ . Vu
les hypothèses faites sur $G$, cet homomorphisme se relève en un homomorphisme
\[ \hand{$\tilde{\rho}$} : G \ \to \ \mathrm{Spin}(q). \]
En utilisant (3.1.2) on déduit de là une application
\[ (3.2.1) \qquad i_{\hand{$\rho$}} : H^1(k, G) \quad \to \ H^3(k, Z/2Z) \ , \]
dont on montre facilement qu'elle ne dépend pas du choix de $q$.

\mspage{\hand{291}}
